\chapter*{Notation and conventions}
\addcontentsline{toc}{chapter}{Notation and conventions}

We collect here the conventions that recur throughout the manual. The complete
symbol list is given in \cref{app:notation}.

\paragraph{Position and velocity.} A particle position is the triple
$\pos=(\lonx,\latx,\depx)$ of longitude, latitude (both in degrees), and depth.
The model velocity is $\velvec=(u,v,\Wvel)$, with horizontal components $u,v$ in
metres per second and vertical component $\Wvel$ in metres per second.

\paragraph{Depth and vertical velocity are positive downward.} Depth $\depx$
increases from $0$ at the surface to the sea floor, and the vertical velocity
$\Wvel$ is positive for downwelling. Consequently the vertical position update
$\depx \mathrel{+}= \Wvel\,\dt$ deepens a particle in downwelling, so that sinking
is followed with the correct sign.

\paragraph{Integration direction.} The factor $\tdir\in\{+1,-1\}$ (the
configuration key \conf{time\_dir}) selects forward ($+1$) or backward ($-1$)
integration. The time step $\dt$ is always a positive quantity; the direction
enters by multiplying the interpolated velocity, so that backward integration is
a sign reversal of the sampled velocity rather than a separate code path
(\cref{ch:arch,ch:integration}).

\paragraph{The spherical metric.} Velocities in metres per second are converted to
angular increments in degrees by the local metric. With
$\Mlat=\SI{111120}{\metre\per\degree}$ the (constant) metres per degree of
latitude and $\zzero=\cos(\latx)\,\Mlat$ the metres per degree of longitude,
the horizontal update reads $\lonx \mathrel{+}= \dt\,u/\zzero$ and
$\latx \mathrel{+}= \dt\,v/\Mlat$. The value of $\Mlat$ follows the nautical-mile
convention ($1852\times 60$), the value used by Parcels; the original Fortran tracker used
\SI{111300}{\metre\per\degree}, so the two differ by \num{0.16}\,\% in every displacement
(\cref{ch:validation}). Diffusive kicks are the exception: they move along great circles
(\cref{sec:diff-sphere}).

\paragraph{The FESOM family.} \textsf{KELP} targets the unstructured-mesh ocean
models of the FESOM family: \emph{FESOM2}, whose meshes are purely triangular and
whose vertical coordinate is a set of fixed depth levels (z-levels); its finite-element predecessor
\emph{FESOM~1.4}, on the same kind of mesh and levels, whose fields live on the levels at the nodes; and
\emph{FESOM-C}, whose meshes are \emph{mixed} --- each element being a triangle or a
quadrilateral --- and whose vertical coordinate is terrain-following (sigma). The
engine additionally reads ocean fields remapped to the HEALPix grid and structured NEMO grids, which it
represents as meshes of quadrilaterals (with a few triangles). Accordingly, throughout this manual an \emph{element} is a
triangle or a quadrilateral, and a ``FESOM mesh'' means a mesh of any of these models;
statements that hold for only one of them name it explicitly. Sizes quoted as examples are those of the
FESOM2 \emph{core2} (triangular, z-level) mesh.

\paragraph{Where fields live.} On a FESOM mesh, scalar fields and node-layout
velocities are defined at the element \emph{vertices} (nodes); element-layout
velocities are defined at element \emph{centroids} (elements); and the vertical
velocity and other interface quantities are defined at the \emph{vertical
interfaces} that bound the depth layers, the three-dimensional horizontal velocity at the layer centres
(FESOM~1.4: on the levels, with the vertical velocity). The distinction governs which
interpolation stencil applies and is made explicit in every schematic.

\paragraph{Typographic conventions.} Source files, functions, and configuration
keys are set in a monospaced font: a configuration key such as \conf{uv\_interp},
a function such as \fn{run\_tracking}, a file such as
\srcfile{numerics/interpolate.py}. A \textbf{Source} callout ties a passage of
prose to the implementation that realizes it.
